\documentclass[11pt]{article}
\usepackage[margin=1in]{geometry}
\usepackage{graphicx,booktabs,float,xcolor,hyperref,amssymb}
\usepackage{helvet}\renewcommand{\familydefault}{\sfdefault}
\hypersetup{colorlinks=true,linkcolor=blue,urlcolor=blue}
\newcommand{\Mw}{$M_{\mathrm w}$}
\setlength{\parindent}{0pt}\setlength{\parskip}{6pt}
\title{\textbf{Source model of the 2026 \Mw~6.9 North Sumatra intraslab earthquake}\\
\large A 172-km-deep rupture beneath the Toba region --- and the cleanest test yet of the
teleseismic MT$+$FFI pipeline}
\author{Automated pipeline (17.Venezuela\_2026)}
\date{2026-08-15 --- no USGS finite-fault product exists for this event}
\begin{document}\maketitle

\section*{Summary}
The \Mw~6.9 earthquake of 2026-08-15 beneath north Sumatra ruptured the subducted
Indo-Australian (Sunda) slab at 172--181 km depth, below the Toba caldera region --- an
intraslab event, not a megathrust rupture. With 231 teleseismic stations at a 21$^\circ$
maximum azimuthal gap and a 98.6\,\% double-couple source, this is the pipeline's cleanest
case: body-wave misfits of 0.080 (lowest of any event studied here), $M_0$ recovered to
$<$1\,\% of the W-phase value, and an \emph{independent} MTUQ mechanism within 7.4$^\circ$
(Kagan) of the W-phase solution --- also the best of any event here. It is the largest deep
earthquake recorded within 150 km of its epicentre (previous maximum \Mw~6.3, 2006).

\section{Event and setting}
\begin{table}[H]\centering\small
\begin{tabular}{ll}
\toprule
Origin & 2026-08-15T10:54:51.8 UTC, 3.086$^\circ$N, 99.015$^\circ$E, depth 172.5 km \\
Magnitude & \Mw~6.9 (W-phase); $M_0=2.62\times10^{19}$ N\,m; centroid 180.5 km; DC 98.6\,\% \\
W-phase NP1 & 207.96/58.15/1.48 \quad (left-lateral, moderately dipping) \\
W-phase NP2 & 117.18/88.75/148.14 \quad (near-vertical, right-lateral oblique) \\
W-phase centroid & 3.080$^\circ$N, 99.013$^\circ$E, origin $+$10.9 s (0.6 km offset) \\
Setting & intraslab, Sunda slab beneath the Toba caldera region \\
\bottomrule
\end{tabular}
\end{table}

The hypocentre--centroid check (Colombia lesson) is benign here: a 0.6 km spatial offset with a
$+$10.9 s time shift is a normal centroid lag for an \Mw~6.9 with a $\sim$7 s half duration,
not a delayed main rupture, so the origin-anchored inversion is valid.

\textbf{Regional significance.} ComCat 1980--2026 within 0.5--5.5$^\circ$N, 96.5--101.5$^\circ$E
lists 718 events \Mw$\geq$4 deeper than 60 km (266 deeper than 150 km): a continuous, active
intermediate-depth zone. Within 150 km of this epicentre and deeper than 120 km the largest
prior event was \Mw~6.3 (2006-12-01, 204 km depth, 35 km away), with all others
\Mw$\leq$5.6 --- this event releases roughly eight times that moment.

\section{Data}
231 three-component stations, 31--90$^\circ$, maximum azimuthal gap 21$^\circ$ --- the best
coverage of any event in this project. The mechanism is a strike-slip pair, and clustered
azimuths are known to alias four-lobed strike-slip radiation (the 2026 Kumamoto lesson); a
21$^\circ$ gap removes that risk. At 180 km the fault lies far below WISP's $\sim$126.5 km
surface-wave Green's-function cap, so surface waves are dropped on \emph{both} planes and the
two-plane comparison is body-only and fair by construction.

\begin{figure}[H]\centering
\includegraphics[width=0.6\textwidth]{../fig_sumatra_stations.png}
\caption{Teleseismic stations (azimuthal-equidistant; rings 30/60/90$^\circ$).}
\end{figure}

\section{Moment tensor (MTUQ)}
\begin{table}[H]\centering\small
\begin{tabular}{llll}
\toprule
variant & mechanism & \Mw & Kagan vs W-phase \\
\midrule
\textbf{DC body P$+$SH @180.5 km} & \textbf{202/56/$-$7} & 6.75 & \textbf{7.4$^\circ$} \\
DC joint @180.5 km & 202/49/2 & 6.75 & 11.3$^\circ$ \\
depth-scan solution at 180.5 km & 207/58/$-$4 & 6.75 & 5.5$^\circ$ \\
\bottomrule
\end{tabular}
\caption{The best independent mechanism recovery in this project. For contrast, the 2026
Colombia \Mw~7.4 gave 52--95$^\circ$ (point source broken by source complexity) and the 2023
Mariana \Mw~6.1 gave 52--86$^\circ$ (signal-to-noise floor). Here a simple 98.6\,\%-DC source,
a 21$^\circ$ azimuthal gap and the raised deep-event body band (25--60 s) combine to recover
the W-phase mechanism almost exactly. \Mw~6.75 versus 6.9 is a modest underestimate.}
\end{table}
\begin{figure}[H]\centering
\includegraphics[width=0.92\textwidth]{../fig_sumatra_beachballs.png}
\caption{W-phase reference and the three MTUQ recoveries.}
\end{figure}

\section{Finite fault (WISP): the fault had to be enlarged}
First-pass \texttt{auto\_model} fits were excellent (misfit 0.080--0.082) but violated the
edge-containment rule: 29\,\% (NP1) and 47\,\% (NP2) of peak slip lay in the deepest fault row.
Following the San Pedro de Atacama rule the fault was enlarged \emph{identically on both
planes} (69$\times$57 km $\to$ 97$\times$114 km) and both were re-inverted and re-shifted.

\begin{table}[H]\centering\small
\begin{tabular}{lll}
\toprule
 & first pass (69$\times$57 km) & enlarged (97$\times$114 km) \\
\midrule
edge slip top / bottom & 2--4\,\% / \textbf{29--47\,\%} & 2--3\,\% / \textbf{2--4\,\%} \\
misfit NP1 / NP2 & 0.0816 / 0.0803 & 0.0829 / \textbf{0.0795} \\
plane margin & 1.6\,\% (meaningless) & \textbf{4.1\,\%, NP2 preferred} \\
peak slip & 0.53--0.57 m & 0.44--0.50 m \\
$\Delta\sigma$ (effective-area) & 1.2 MPa & 0.65--0.83 MPa \\
\bottomrule
\end{tabular}
\caption{The enlargement changed the conclusion: with the boundary artefact removed the plane
margin grew from a meaningless tie to a suggestive 4.1\,\% favouring NP2 (208/58/1) --- the
plane MTUQ independently converged on. Both models recover $M_0=2.64\times10^{19}$ N\,m
($<$1\,\% from the W-phase) with T$_{5-95}\approx$15.5 s, and the source-time function is
plane-independent ($r=0.972$): well-determined moment release on a weakly-determined plane.}
\end{table}

\begin{figure}[H]\centering
\includegraphics[width=0.82\textwidth]{../20260815105451/ffm.0/NP2_big_shift/plots/SlipIsochrone_trueScale.png}
\caption{Preferred plane (208/58/1), true scale: compact slip at 165--207 km depth, peak
0.50 m, cleanly contained after enlargement (edge slip 2\,\%).}
\end{figure}
\begin{figure}[H]\centering
\includegraphics[width=0.55\textwidth]{../20260815105451/ffm.0/NP2_big_shift/plots/MomentRate.png}
\caption{Moment-rate function (preferred plane).}
\end{figure}
\begin{figure}[H]\centering
\includegraphics[width=0.92\textwidth]{../20260815105451/ffm.0/NP2_big_shift/plots/P_body_waves.png}
\caption{Teleseismic P fits (preferred plane, shifted).}
\end{figure}
\begin{figure}[H]\centering
\includegraphics[width=0.92\textwidth]{../20260815105451/ffm.0/NP2_big_shift/plots/SH_body_waves.png}
\caption{Teleseismic SH fits (preferred plane, shifted).}
\end{figure}
\begin{figure}[H]\centering
\includegraphics[width=0.82\textwidth]{../20260815105451/ffm.0/NP1_big_shift/plots/SlipIsochrone_trueScale.png}
\caption{Conjugate plane (117/89/148) for comparison: 4.1\,\% worse fit.}
\end{figure}

\section{Assessment}
\begin{figure}[H]\centering
\includegraphics[width=\textwidth]{../fig_sumatra_assessment.png}
\caption{Left: the depth scan --- misfit (blue) is flat to 0.3\,\% across 140--220 km, but the
Kagan angle (red) reveals a sharp transition: the mechanism is chaotic at 140--170 km
(85--99$^\circ$) and locks onto the W-phase from 180.5 km downward (5.5--18$^\circ$). Centre:
the source-time function is plane-independent. Right: the plane-discrimination margin in the
project's rupture-size context.}
\end{figure}

\textbf{Depth --- a split verdict worth stating carefully.} The body-wave misfit minimum does
\emph{not} resolve depth (0.3\,\% variation over an 80-km range; its nominal minimum at 155 km
is meaningless at that flatness). What does resolve it is \emph{mechanism stability}: the body
waves only fit coherently once the depth phases are correctly timed, which occurs at
$\geq$180 km. Together with the WISP slip centroid at 181 km, this independently corroborates
the USGS 180.5 km centroid over the 172.5 km hypocentre.

\textbf{Kinematics (preferred plane).} The slip centroid lies 8.4 km from the epicentre at
azimuth 349$^\circ$, 181 km deep; the moment centroid migrates 13.7 km toward azimuth
354$^\circ$ (north) while deepening from 176 to 188 km, with 66\,\%/23\,\% forward/backward
moment --- a moderately unilateral, northward and down-dip rupture. The realized rupture-front
speed is 2.52 km s$^{-1}$ (16--84\,\%: 2.24--2.95), giving $V_r/V_s\approx0.55$ at
$V_s\approx4.6$ km s$^{-1}$: the same slow, dissipative regime found for the 2026 Colombia
\Mw~7.4 (0.56), and matching the slow first stage rather than the fast runaway stage of
Poli et al.\ (2016). The stress drop, 0.65--0.83 MPa, is low.

\section{Source-time function, asperities, stress drop, rupture speed}
\textbf{The double-peaked STF is robust; two comparable asperities are not.} The second
moment-rate pulse appears in four of five independently derived models (both planes $\times$
both fault sizes) and its prominence \emph{grows} as the parameterization improves (0.07 and
0.16 of peak in the original edge-violating faults; 0.24 and 0.38 after enlargement). Splitting
the moment by rupture time, the first pulse (t$_{rup}<$11 s) carries 67\,\% of the moment at
165--194 km depth and the second (t$_{rup}\geq$11 s) carries 33\,\% at 191--207 km; their
centroids are 20 km apart (15 km horizontal, 13 km deeper) separated by 9.3 s, an apparent
migration of 2.15 km s$^{-1}$ --- a legitimate rupture velocity, in contrast to the
kinematically impossible 0.8 km s$^{-1}$ of the 2026 Colombia \Mw~7.4. As a \emph{spatially
distinct} asperity, however, the second patch is marginal: at a 50\,\%-of-peak threshold it is
a single subfault and at 60\,\% it disappears. The defensible statement is a dominant asperity
plus a weaker, deeper secondary pulse $\sim$20 km down-dip.

\textbf{Stress drop depends on convention by an order of magnitude.}

\begin{table}[H]\centering\small
\begin{tabular}{llll}
\toprule
convention & area & radius & $\Delta\sigma$ \\
\midrule
effective, $(\Sigma u)^2/\Sigma u^2$ & 1827 km$^2$ & 24.1 km & \textbf{0.82 MPa} \\
slip $>$15\,\% of peak & 1479 km$^2$ & 21.7 km & 1.13 MPa \\
half-peak ($>$50\,\%) & 440 km$^2$ & 11.8 km & 6.98 MPa \\
\bottomrule
\end{tabular}
\caption{The spread follows from smooth slip (0.20 m mean over the effective area versus a
0.50 m peak). A slip-weighted average of the on-fault static stress-change field (Ripperger \&
Mai 2007 kernel, Noda et al.\ 2013 energy-based averaging) gives $\approx$0.9 MPa for the
preferred model: the effective-area convention tracks the rigorous value, while the half-peak
proxy overestimates it by nearly an order of magnitude. All are lower bounds, since inversion
smoothing suppresses short-wavelength slip and hence stress.}
\end{table}

\textbf{Rupture speed} is stable across models and interior to its allowed window
(1.2--3.8 km s$^{-1}$, reference 3.0), so it is data-selected rather than a bound artefact:
2.52 (preferred NP2; 16--84\,\% 2.24--2.95), 2.60 (NP1), 2.50 and 2.38 km s$^{-1}$ on the
smaller faults. With $V_s\approx$4.6 km s$^{-1}$ at 180 km this gives
$V_r/V_s\approx$\textbf{0.55} --- the same slow, dissipative regime as the Colombia \Mw~7.4
(0.56).

\begin{figure}[H]\centering
\includegraphics[width=0.42\textwidth]{../../results/mtuq_dc_2026sumatra/2026sumatra_waveforms_body.png}
\caption{Waveform fits of the adopted MTUQ solution (202/56/$-$7, body P$+$SH), rendered at the
fixed adopted solution rather than whichever variant ran last. Black observed, red synthetic.}
\end{figure}

\section{Caveats and next steps}
Teleseismic-only; no USGS finite fault exists for comparison; the plane preference (4.1\,\%) is
suggestive rather than decisive, though corroborated by the independent MTUQ mechanism; depth
rests on mechanism stability rather than a misfit minimum; peak slip is smoothing-limited and
therefore a lower bound. Natural next steps: a radiated-energy and efficiency measurement (does
this event also sit in the low-$\eta$ dissipative regime found for Colombia?); BMKG regional
waveforms for an independent plane test; and comparison with the 2006 \Mw~6.3 mechanism to test
whether this slab segment ruptures repeatedly in the same geometry.

\end{document}
